\documentclass{article}
\usepackage[margin=1in]{geometry}
\usepackage[utf8]{inputenc}
\usepackage{amsmath}
\usepackage{graphicx}

\title{ELCT 432 PA9}
\author{Jacob Vaught}
\date{12-03-2023}

\begin{document}

\maketitle

\subsection*{Question 1: Coin Tossing}
\textbf{a. Amount of Information for "head" and "tail":}

The amount of information for an outcome in a random process is given by the formula:
\[ I(x) = -\log_2 P(x) \]
where \( P(x) \) is the probability of occurrence of outcome \( x \).

For a coin toss where the probability of getting a head is 0.2:
\begin{align*}
    I(\text{head}) &= -\log_2(0.2) \approx 2.322 \text{ bits} \\
    I(\text{tail}) &= -\log_2(0.8) \approx 0.322 \text{ bits}
\end{align*}

\textbf{b. Average Information (Entropy):}
The average information (entropy) of a random process is given by:
\[ H(X) = -\sum P(x) \log_2 P(x) \]
where the sum is over all possible outcomes \( x \).

For the coin toss:
\begin{align*}
    H(X) &= -(0.2 \log_2 0.2 + 0.8 \log_2 0.8) \\
    &= -(-2.322 \times 0.2 - 0.322 \times 0.8) \\
    &= 0.722 \text{ bits}
\end{align*}
\subsection*{Question 2: Binary Symmetric Channel (BSC)}
Given a BSC with \( P(0|1) = P(1|0) = \epsilon = 0.3 \) and \( P(X=1) = P(X=0) = 0.5 \).

\textbf{a. Amount of Information in \( Y \):}
First, calculate \( P(Y=0) \) and \( P(Y=1) \), then use these to find the entropy \( H(Y) \).

The probabilities can be calculated as:
\begin{align*}
    P(Y=0) &= P(X=0)(1 - \epsilon) + P(X=1)\epsilon = 0.5 \times 0.7 + 0.5 \times 0.3 = 0.5 \\
    P(Y=1) &= 1 - P(Y=0) = 0.5
\end{align*}

The entropy \( H(Y) \) is then given by:
\begin{align*}
    H(Y) &= -P(Y=0)\log_2(P(Y=0)) - P(Y=1)\log_2(P(Y=1)) \\
    &= -0.5 \log_2(0.5) - 0.5 \log_2(0.5) = 1 \text{ bit}
\end{align*}

\textbf{b. Amount of Information in \( X \) Given \( Y \):}
The conditional entropy \( H(X|Y) \) is calculated using the joint probabilities and marginal probabilities.

Calculating joint probabilities:
\begin{align*}
    P(X=0, Y=0) &= P(X=0)(1 - \epsilon) = 0.35 \\
    P(X=0, Y=1) &= P(X=0)\epsilon = 0.15 \\
    P(X=1, Y=0) &= P(X=1)\epsilon = 0.15 \\
    P(X=1, Y=1) &= P(X=1)(1 - \epsilon) = 0.35
\end{align*}

Then, the joint entropy \( H(X, Y) \) is:
\begin{align*}
    H(X, Y) &= -\sum P(X=x, Y=y) \log_2 P(X=x, Y=y) \\
    &= -0.35 \log_2(0.35) - 0.15 \log_2(0.15) \\
    &\quad - 0.15 \log_2(0.15) - 0.35 \log_2(0.35) \approx 1.881 \text{ bits}
\end{align*}

The conditional entropy \( H(X|Y) \) is:
\begin{align*}
    H(X|Y) &= H(X, Y) - H(Y) = 1.881 - 1 = 0.881 \text{ bits}
\end{align*}

\textbf{c. Uncertainty Reduction on \( X \) After Observing \( Y \):}
The mutual information \( I(X; Y) \) is calculated as:
\begin{align*}
    I(X; Y) &= H(X) - H(X|Y) \\
    &= 1 - 0.881 = 0.119 \text{ bits}
\end{align*}

The mutual information represents the amount of uncertainty reduced in \( X \) after observing \( Y \).
\subsection*{Question 3: Uniform Quantization}
Uniform quantization is not a linear operation. In linear operations, the output is directly proportional to the input. However, in uniform quantization, a continuous range of input values is mapped to discrete levels. This mapping results in a loss of detailed information and introduces quantization error, especially noticeable when the input varies smoothly within a quantization range. Hence, uniform quantization is inherently nonlinear.

\subsection*{Question 4: Channel Capacity of a Gaussian Channel}
For a Gaussian channel where \( Y = X + Z \), and \( Z \) is a zero-mean Gaussian random variable with variance \( \sigma_n^2 \), the channel capacity can be derived using the Shannon-Hartley theorem.

Given the power constraint \( \frac{1}{n} \sum_{i=1}^{n} x_i^2 \leq P \), the channel capacity \( C \) is given by the formula:
\[ C = \frac{1}{2} \log_2 \left(1 + \frac{P}{\sigma_n^2}\right) \]
This formula is derived by maximizing the mutual information between \( X \) and \( Y \), considering Gaussian distributions for the noise and signal. The proof involves calculating the differential entropy of the Gaussian noise and using the properties of logarithmic functions. 

\subsection*{Question 5: Binary Symmetric Channel Analysis}
Given a BSC with \( P(X=1) = 0.8, P(X=0) = 0.2 \), and bit flip probability 0.25.

\textbf{a. Amount of Information for Observing "1" at the Source:}
\begin{align*}
    I(1) &= -\log_2(P(X=1)) = -\log_2(0.8) \approx 0.322 \text{ bits}
\end{align*}

\textbf{b. Amount of Information for Observing "0" at the Source:}
\begin{align*}
    I(0) &= -\log_2(P(X=0)) = -\log_2(0.2) \approx 2.322 \text{ bits}
\end{align*}

\textbf{c. Average Information at the Source \( H(X) \):}
\begin{align*}
    H(X) &= -P(X=0)\log_2(P(X=0)) - P(X=1)\log_2(P(X=1)) \\
    &= -(0.2 \log_2 0.2 + 0.8 \log_2 0.8) \\
    &= 0.722 \text{ bits}
\end{align*}

\textbf{d, e. Probability that the Output of the Channel is "1" and "0":}
\begin{align*}
    P(Y=1) &= P(X=1)(1 - \text{flip probability}) + P(X=0)(\text{flip probability}) \\
    &= 0.8 \times 0.75 + 0.2 \times 0.25 = 0.65 \\
    P(Y=0) &= 1 - P(Y=1) = 0.35
\end{align*}

\textbf{f. Average Information at the Output \( H(Y) \):}
\begin{align*}
    H(Y) &= -P(Y=0)\log_2(P(Y=0)) - P(Y=1)\log_2(P(Y=1)) \\
    &= -(0.35 \log_2 0.35 + 0.65 \log_2 0.65) \\
    &= 0.934 \text{ bits}
\end{align*}
\textbf{g. Probability of Transmitting 0 and Receiving 0:}
\begin{align*}
    P(X=0, Y=0) &= P(X=0)(1 - \text{flip probability}) \\
    &= 0.2 \times 0.75 = 0.15
\end{align*}

\textbf{h. Probability of Transmitting 1 and Receiving 0:}
\begin{align*}
    P(X=1, Y=0) &= P(X=1)(\text{flip probability}) \\
    &= 0.8 \times 0.25 = 0.2
\end{align*}

\textbf{i. Probability of Transmitting 0 and Receiving 1:}
\begin{align*}
    P(X=0, Y=1) &= P(X=0)(\text{flip probability}) \\
    &= 0.2 \times 0.25 = 0.05
\end{align*}

\textbf{j. Probability of Transmitting 1 and Receiving 1:}
\begin{align*}
    P(X=1, Y=1) &= P(X=1)(1 - \text{flip probability}) \\
    &= 0.8 \times 0.75 = 0.6
\end{align*}

\textbf{k. Joint Entropy \( H(X, Y) \):}
The joint entropy \( H(X, Y) \) is calculated using the joint probabilities:
\begin{align*}
    H(X, Y) &= -\sum P(X=x, Y=y) \log_2 P(X=x, Y=y) \\
    &= -(P(X=0, Y=0) \log_2 P(X=0, Y=0) + P(X=0, Y=1) \log_2 P(X=0, Y=1) \\
    &\quad + P(X=1, Y=0) \log_2 P(X=1, Y=0) + P(X=1, Y=1) \log_2 P(X=1, Y=1)) \\
    &= -(0.15 \log_2 0.15 + 0.05 \log_2 0.05 + 0.2 \log_2 0.2 + 0.6 \log_2 0.6) \\
    &\approx 1.533 \text{ bits}
\end{align*}

\textbf{l. Mutual Information \( I(X; Y) \):}
The mutual information \( I(X; Y) \) is calculated as the difference between the entropy of \( X \) and the conditional entropy of \( X \) given \( Y \):
\begin{align*}
    I(X; Y) &= H(X) - H(X|Y) \\
    &= H(X) - (H(X, Y) - H(Y)) \\
    &= 0.722 - (1.533 - 0.934) \\
    &= 0.123 \text{ bits}
\end{align*}

This mutual information represents the amount of uncertainty about \( X \) that is removed by observing \( Y \) in the binary symmetric channel.

\end{document}
